\section{How to Use Our Code}

% Include important code snippets such that we know how to use your code.

\subsection{System Architecture and Overview}
\label{sec:use-architecture}

% Overview over directories and files in the repository: file structure
% Class diagram for different modules: showing realization of abstract quantum\_interpreter etc.

\subsection{Environment Setup}
\label{sec:use-env-setup}

To ensure dependencies do not conflict with your system-wide Python installation, we recommend executing this project within an isolated virtual environment. 
From the project's root directory, initialize and activate the environment, then install the package along with its required dependencies in editable mode by running:

\begin{lstlisting}[
    language=bash, 
    numbers=none, 
    backgroundcolor={},
    caption={Terminal installation steps}, 
    label={lst:install_bash}
]
# 1. Create a virtual environment named .venv
python -m venv .venv

# 2. Activate the environment
# On Linux / macOS:
source .venv/bin/activate
# On Windows (PowerShell):
.venv\Scripts\Activate.ps1

# 3. Install the project dynamically from pyproject.toml
pip install -e .
\end{lstlisting}

By utilizing the \texttt{-e} flag, the source code inside the \texttt{src/} directory is linked directly to your active environment. This allows any absolute imports to resolve automatically from anywhere on your system.

\subsection{Using the Interpreter to create a GHZ State}
\label{sec:use-interpreter}

In the following we will look at the script \texttt{examples/ghz.py} (see Listing~\ref{lst:ghz_example}) to see how to interpret a brainfuq program that creates a GHZ-State and how to translate it into a qiskit circuit.

\begin{lstlisting}[language=Python, caption=examples/ghz.py, label={lst:ghz_example}]
from brainfuq_interpreter.brainfuq_interpreter import BrainfuqInterpreter
from brainfuq_interpreter.brainfuq_simulator import BrainfuqSimulator
from brainfuq_interpreter.brainfuq_to_qiskit_translator import QiskitBrainfuqInterpreter

N = 10

draw = True
output_method = "latex_source" # default console output for None or "text"

simulator = BrainfuqSimulator()
interpreter = BrainfuqInterpreter(simulator)

# put q0 in equal superposition
program = "~"

# loop counter
program += "+" * (N - 1)

# loop body: add a cnot and decrement loop counter
program += "[#}*-]"

interpreter.interpret_program(program)
print(simulator)

# reset interpreter before interpreting the next program
interpreter.reset()

if draw:
    translation_interpreter = BrainfuqInterpreter(QiskitBrainfuqInterpreter(verbose=False))
    qc, _ = translation_interpreter.interpret_program(program)
    print(qc.draw(output=output_method))

\end{lstlisting}

Running the script outputs \texttt{0.707 |0000000000⟩ + 0.707 |1111111111⟩} as well as qiskit circuit printed as latex source code (s. \Cref{fig:ghz_circuit}).

\begin{figure}[htbp]
\centering
\scalebox{1.0}{
\Qcircuit @C=1.0em @R=0.2em @!R { \\
	 	\nghost{{q}_{0} :  } & \lstick{{q}_{0} :  } & \gate{\mathrm{H}} & \ctrl{1} & \qw & \qw & \qw & \qw & \qw & \qw & \qw & \qw & \qw & \qw\\
	 	\nghost{{q}_{1} :  } & \lstick{{q}_{1} :  } & \qw & \targ & \ctrl{1} & \qw & \qw & \qw & \qw & \qw & \qw & \qw & \qw & \qw\\
	 	\nghost{{q}_{2} :  } & \lstick{{q}_{2} :  } & \qw & \qw & \targ & \ctrl{1} & \qw & \qw & \qw & \qw & \qw & \qw & \qw & \qw\\
	 	\nghost{{q}_{3} :  } & \lstick{{q}_{3} :  } & \qw & \qw & \qw & \targ & \ctrl{1} & \qw & \qw & \qw & \qw & \qw & \qw & \qw\\
	 	\nghost{{q}_{4} :  } & \lstick{{q}_{4} :  } & \qw & \qw & \qw & \qw & \targ & \ctrl{1} & \qw & \qw & \qw & \qw & \qw & \qw\\
	 	\nghost{{q}_{5} :  } & \lstick{{q}_{5} :  } & \qw & \qw & \qw & \qw & \qw & \targ & \ctrl{1} & \qw & \qw & \qw & \qw & \qw\\
	 	\nghost{{q}_{6} :  } & \lstick{{q}_{6} :  } & \qw & \qw & \qw & \qw & \qw & \qw & \targ & \ctrl{1} & \qw & \qw & \qw & \qw\\
	 	\nghost{{q}_{7} :  } & \lstick{{q}_{7} :  } & \qw & \qw & \qw & \qw & \qw & \qw & \qw & \targ & \ctrl{1} & \qw & \qw & \qw\\
	 	\nghost{{q}_{8} :  } & \lstick{{q}_{8} :  } & \qw & \qw & \qw & \qw & \qw & \qw & \qw & \qw & \targ & \ctrl{1} & \qw & \qw\\
	 	\nghost{{q}_{9} :  } & \lstick{{q}_{9} :  } & \qw & \qw & \qw & \qw & \qw & \qw & \qw & \qw & \qw & \targ & \qw & \qw\\
\\ }}
\caption{Rendering of the latex source code produced for the qiskit circuit translated from brainfuq}
\label{fig:ghz_circuit}
\end{figure}

%\subsection{Translating between Brainfuq and Qiskit}
%\label{sec:use-translator}

